\documentclass[../main.tex]{subfiles}
\graphicspath{{\subfix{../images/}}}
\begin{document}

  Os artigos deverão ser enviados em formato digitalizado (Word) com texto  completo.  O  idioma  para  artigos  submetidos  é  o  Português.  Nenhuma  revisão  dos  textos  será  feita,  sendo  os  mesmos  de  inteira  responsabilidade  de  seus  autores.
      
  O limite de páginas para o artigo está fixado entre 6 (seis) e 15 (quinze)  para defesa junto a banca acadêmica, considerando-se texto, tabelas, figuras e  fotos.  Estas  páginas  deverão  estar  configuradas  para  formato  A4  (210  x  297  mm), com margens laterais e superior de 25mm e inferior de 30mm. Não deverá  haver nelas nenhum texto nas áreas de cabeçalho e rodapé. As figuras poderão  ser  coloridas  e  deverão  ser  inseridas  no  corpo  do  trabalho,  próximas  as  suas  citações no texto. As figuras deverão ser centralizadas, sem exceder o tamanho  limitado pelas margens da página.

  Cada  figura  deverá  ter  um  título  numerado  em  algarismos  arábicos.  Os  títulos deverão ser escritos em fonte Times New Roman / Arial 11, itálico e devem  ser centralizados na parte inferior da figura. 

  \begin{figure}[h]
    \centering
    \caption{Titulo da Figura}
    \includegraphics[width=0.5\textwidth]{logo_cimatec.png}
    
    Fonte: Proprio autor
    \label{fig:fig1}
  \end{figure}

  As tabelas deverão ser justificadas (usando toda a área entre as  margens).  As  tabelas  poderão  ser  coloridas  ou  em  preto  e  branco.  O  título  da  tabela  deve  ser  alinhado  à  margem  esquerda  e  digitado  em  fonte  Times  New  Roman  /  Arial  11.  As  unidades  de  medida  correspondentes  a  todos  os  termos  deverão ser claramente indicadas, usando o sistema internacional (S.I.). 

  \begin{table}
    \caption{Titulo da Tabela}
    \centering
    \begin{tabular}{ |c|c|c|c| } 
      \hline
      col1 & col2 & col3 \\
      \hline
      \multirow{3}{4em}{Multiple row} & cell2 & cell3 \\ 
      & cell5 & cell6 \\ 
      & cell8 & cell9 \\ 
      \hline
      \end{tabular}
    
      Fonte: Proprio autor
    \label{tab:table-name}
  \end{table}


\end{document}